% ============================================================
% 第2章第7節「特殊関数型」草案
% ============================================================
% このファイルは現行本文には接続していない草案である.
% 特殊関数の理論的な展開は付録へ移し,ここでは漸化式を解くための
% 具体的な見抜き方と,候補を帰納法で確定する流れに集中する.

\section{特殊関数型}

ここまで,漸化式の形を見て,等差型,等比型,階差型,階比型などの
基本パターンへ帰着する方法を学んできた.しかし,どのパターンにも
そのまま入らない漸化式もある.

例えば,
\[
  a_{n+1}=a_n^2
\]
ならば,対数を取って等比型にできる.一方,
\[
  a_{n+1}=a_n^2+1
\]
では,対数を取っても定数項が残る.このようなときは,式を無理に
基本パターンへ押し込むのではなく,右辺と一致する恒等式を探す.

\begin{mdquote}
特殊関数型では,まず「次の項を作る右辺」がどの公式と一致するかを
探す.次に,その公式に合わせて数列を表す候補を作り,初項と漸化式を
帰納法で確認する.公式を思い出すだけでは解答にならない.
\end{mdquote}

\textbf{パターン1: 倍角公式が見えたら,初項を余弦に直す}

次の漸化式を解く.
\[
  a_1=\frac{\sqrt{3}}{2},
  \qquad
  a_{n+1}=2a_n^2-1.
  \tag{2.7.1}
\]
右辺を見たら,倍角公式
\[
  \cos 2\theta=2\cos^2\theta-1
  \tag{2.7.2}
\]
を思い出す.初項も
\[
  a_1=\cos\frac{\pi}{6}
\]
と書けるので,
\[
  a_n=\cos\left(2^{n-1}\frac{\pi}{6}\right)
  \tag{2.7.3}
\]
を候補にする.

候補が正しいことを帰納法で示す.\,$n=1$ では初項の式そのものである.
ある $n$ で $a_n=\cos(2^{n-1}\pi/6)$ が成り立つとすると,
\begin{align*}
  a_{n+1}
    &=2a_n^2-1\\
    &=2\cos^2\left(2^{n-1}\frac{\pi}{6}\right)-1\\
    &=\cos\left(2^n\frac{\pi}{6}\right).
\end{align*}
これは $n+1$ の式である.したがって,

\begin{mdquote}
\[
  a_n=\cos\left(2^{n-1}\frac{\pi}{6}\right)
\]
\end{mdquote}

となる.ここで行ったことは,次の4段階にまとめられる.

\begin{enumerate}
  \item 右辺 $2x^2-1$ を見て,倍角公式を思い出す.
  \item 初項を $\cos\theta_1$ の形に直す.
  \item 角度が $\theta_{n+1}=2\theta_n$ と進むと予想する.
  \item 倍角公式を使って,候補を帰納法で確認する.
\end{enumerate}

初項が一般に $a_1=\cos\theta$ なら,同じ計算から
\[
  a_n=\cos(2^{n-1}\theta)
  \tag{2.7.4}
\]
を得る.\,$-1\leq a_1\leq1$ なら,実数 $\theta$ を選んで
$a_1=\cos\theta$ とできる.\,$a_1>1$ の場合は
\[
  \cosh 2t=2\cosh^2t-1
\]
を使い,\,$a_1=\cosh t$ と置けば
$a_n=\cosh(2^{n-1}t)$ となる.初項の範囲を確認してから,置換を選ぶことが大切である.

\textbf{パターン2: 三倍角公式が見えたら,角度を3倍する}

次の形を考える.
\[
  a_1=\cos\theta,
  \qquad
  a_{n+1}=4a_n^3-3a_n.
  \tag{2.7.5}
\]
三倍角公式
\[
  \cos 3\theta=4\cos^3\theta-3\cos\theta
  \tag{2.7.6}
\]
が右辺と一致している.したがって,
\[
  a_n=\cos(3^{n-1}\theta)
  \tag{2.7.7}
\]
と予想する.

実際,\,$n=1$ では初項に一致する.\,$a_n=\cos(3^{n-1}\theta)$ と仮定すると,
\begin{align*}
  a_{n+1}
    &=4a_n^3-3a_n\\
    &=4\cos^3(3^{n-1}\theta)-3\cos(3^{n-1}\theta)\\
    &=\cos(3^n\theta).
\end{align*}
よって候補は漸化式を満たす.

倍角と三倍角の違いは,公式を覚えているかどうかではない.右辺を
$F(x)$ と書いたとき,\,$F(\cos\theta)=\cos(m\theta)$ となる $m$ を
見つければ,角度の漸化式は $\theta_{n+1}=m\theta_n$ になる.
この形をさらに一般化したものが付録のチェビシェフ多項式である.

\textbf{パターン3: 三項漸化式なら,余弦と正弦を組み合わせる}

次のように前の2項が現れる漸化式を考える.
\[
  a_{n+1}=2x a_n-a_{n-1}.
  \tag{2.7.8}
\]
$|x|<1$ なら $x=\cos\theta$ と置ける.加法定理から,
\[
  \cos((n+1)\theta)
    =2\cos\theta\cos(n\theta)-\cos((n-1)\theta)
\]
および
\[
  \sin((n+1)\theta)
    =2\cos\theta\sin(n\theta)-\sin((n-1)\theta)
\]
が成り立つ.したがって,
\[
  a_n=A\cos(n\theta)+B\sin(n\theta)
  \tag{2.7.9}
\]
と置けば,漸化式そのものは自動的に満たされる.

残るのは初期条件から $A,B$ を決めることである.\,$a_0,a_1$ が与えられ,
$\sin\theta\neq0$ なら,
\[
  A=a_0,
  \qquad
  B=\frac{a_1-a_0\cos\theta}{\sin\theta}
  \tag{2.7.10}
\]
である.初期条件を代入して定数を決める,という線形漸化式の基本方針が
ここでも働いている.

この形で重要なのは,特定の初項に合わせて $\cos(n\theta)$ だけを置かない
ことである.一般の初期条件には,余弦と正弦の2つを組み合わせる必要がある.
$\sin\theta=0$ の場合は,特性方程式または $x=1,-1$ の場合分けに戻る.

\textbf{パターン4: $a_n^2-2$ が見えたら,対称な指数を置く}

三角関数とは別に,次の恒等式がある.
\[
  (u+u^{-1})^2-2=u^2+u^{-2}.
  \tag{2.7.11}
\]
そこで,
\[
  a_1=u+u^{-1},
  \qquad
  a_{n+1}=a_n^2-2
  \tag{2.7.12}
\]
に対して,
\[
  a_n=u^{2^{n-1}}+u^{-2^{n-1}}
  \tag{2.7.13}
\]
と置く.\,$n=1$ は初項に一致し,帰納法の段階では
\begin{align*}
  a_{n+1}
    &=\left(u^{2^{n-1}}+u^{-2^{n-1}}\right)^2-2\\
    &=u^{2^n}+u^{-2^n}
\end{align*}
となる.指数が2倍される構造が,そのまま角度を2倍する方法に対応している.

同様に,
\[
  (u+u^{-1})^3-3(u+u^{-1})=u^3+u^{-3}
\]
なので,
\[
  a_{n+1}=a_n^3-3a_n
\]
なら指数を $3$ 倍する.右辺の係数を見て,対称な指数式を保つために
必要な項を逆算するのがこのパターンの方針である.

\textbf{パターン5: 一般項が必要なければ,近づく数列を設計する}

方程式 $x^2=A$ の解を求めたいとする.\,$a_n$ を近似値とし,
接線が $x$ 軸と交わる点を次の近似値に選ぶと,
\[
  a_{n+1}=\frac12\left(a_n+\frac{A}{a_n}\right)
  \tag{2.7.14}
\]
を得る.例えば $A=2,a_1=2$ なら,
\[
  a_2=\frac32,
  \qquad
  a_3=\frac{17}{12},
  \qquad
  a_4=\frac{577}{408}
\]
となり,\,$\sqrt2$ に近づく.

この方法では,一般項を閉じた式で書くことよりも,次の項をどう設計すれば
誤差が小さくなるかを考えている.接線と収束の詳しい議論は付録で扱う.

\textbf{特殊関数型を見たときの最終確認}

特殊関数型の解法は,次の順に進める.

\begin{enumerate}
  \item 右辺を見て,倍角公式,三倍角公式,対称式などとの一致を探す.
  \item 初項がその置換の範囲に入っているかを確認する.
  \item $a_n=\varphi(t_n)$ などの候補を作る.
  \item 公式を使って $t_n$ の漸化式を解く.
  \item 初項と漸化式を帰納法で確認する.
\end{enumerate}

「似ている」だけでは不十分である.実際に恒等式が一致し,候補が初項と
漸化式の両方を満たすことまで示して,初めて解答になる.特殊関数の名前を
知っていることよりも,どの表現なら次の項を同じ形で作れるかを見抜くことが,
このパターンの中心である.
